\documentclass[12pt,a4paper]{article}
\usepackage{fontspec}
\setmainfont{Times New Roman}
\setmonofont{Consolas}[Scale=0.8]
\usepackage[russian]{babel}
\usepackage[left=30mm,right=15mm,top=20mm,bottom=20mm]{geometry}
\usepackage{setspace}\onehalfspacing
\usepackage{tikz}\usetikzlibrary{arrows.meta,positioning}
\usepackage{listings,xcolor,float,hyperref}
\hypersetup{hidelinks,breaklinks=true}
\setlength{\parindent}{1.25cm}
\emergencystretch=2em
\definecolor{codegray}{RGB}{247,247,247}
\lstset{language=C,basicstyle=\ttfamily\footnotesize\singlespacing,backgroundcolor=\color{codegray},breaklines=true,showstringspaces=false,frame=single,framesep=5pt,columns=fullflexible,keepspaces=true}
\begin{document}
\begin{titlepage}
\begin{center}
{\small МИНИСТЕРСТВО НАУКИ И ВЫСШЕГО ОБРАЗОВАНИЯ\\РОССИЙСКОЙ ФЕДЕРАЦИИ}\\[4mm]
{\small Федеральное государственное автономное образовательное учреждение\\ высшего образования\\ «Московский авиационный институт (национальный исследовательский университет)»}
\vfill
{\bfseries\large ОТЧЁТ ПО ЛАБОРАТОРНОЙ РАБОТЕ № 4}\\[9mm]
Дисциплина:\\[2mm]
{\bfseries Операционные системы и архитектура компьютера}\\[6mm]
Тема:\\[2mm]
{\bfseries Межпроцессное взаимодействие с использованием разделяемой памяти}\\[3mm]
Вариант 19
\vfill
\end{center}
\begin{flushright}\begin{minipage}{0.45\textwidth}
Выполнил:\\ студент группы М80-203БВ-25\\ Прудников Я. А.\par\vspace{7mm}
Преподаватель:\\ Миронов Е. С.
\end{minipage}\end{flushright}
\vfill\begin{center}Москва, 2026 г.\end{center}
\end{titlepage}
\section{Цель работы}
Изучить создание процессов и обмен данными между ними через отображаемую разделяемую память. Разработать программу на языке C, в которой родительский процесс распределяет строки между двумя дочерними процессами, а дочерние процессы обрабатывают строки и записывают результат в файлы.
\section{Постановка задачи}
Вариант 19 сохраняет условие лабораторной работы № 2: родительский процесс получает два имени выходных файлов, создаёт два дочерних процесса и передаёт каждому из них строки. Каждая строка с вероятностью 80\% направляется первому процессу и в остальных случаях -- второму. Дочерние процессы удаляют из строк гласные и записывают результат в соответствующий файл. В работе № 4 вместо каналов используется отображаемая память; системные ошибки необходимо обрабатывать.
\section{Архитектура решения}
В проекте родительская и дочерняя программы находятся в файлах \texttt{parent.c} и \texttt{child.c}. Общие функции чтения и обработки строк вынесены в \texttt{line\_utils.c}, а интерфейс работы с процессами и памятью -- в \texttt{os\_api.h}. Реализация для Linux находится в \texttt{platform/os\_linux.c}. Исходный код проекта: \url{https://github.com/yalama34/os-lab-3}.

Родитель считывает имена файлов, проверяет, что они непустые и различаются, затем создаёт два объекта POSIX shared memory. Для каждого объекта выделяется собственный отображённый буфер. Через \texttt{fork()} создаётся процесс, в нём стандартный вывод перенаправляется в файл вызовом \texttt{dup2()}, после чего запускается программа \texttt{child} через \texttt{execvp()}. Имя объекта памяти передаётся ей аргументом. Родитель читает строки до конца ввода и выбирает адресата с помощью \texttt{rand() \% 100}: значения от 0 до 79 соответствуют первому процессу.

Дочерний процесс открывает объект памяти, собирает переданную строку и удаляет символы \texttt{a, e, i, o, u, y} и их прописные варианты. Затем он записывает результат в стандартный вывод, который уже направлен в файл. Обработка кириллических гласных исходным кодом не предусмотрена.
\subsection{Схема взаимодействия компонентов}
\begin{figure}[H]\centering
\begin{tikzpicture}[>=Stealth,node distance=8mm and 13mm, every node/.style={font=\small}, box/.style={draw,rounded corners,align=center,minimum width=36mm,minimum height=9mm,fill=black!3}]
\node[box] (p) {Родитель\\ввод и выбор адресата};
\node[box,below left=of p] (m1) {Память 1\\empty / full};
\node[box,below right=of p] (m2) {Память 2\\empty / full};
\node[box,below=of m1] (c1) {Дочерний процесс 1};
\node[box,below=of m2] (c2) {Дочерний процесс 2};
\node[box,below=of c1] (f1) {Файл 1};
\node[box,below=of c2] (f2) {Файл 2};
\draw[->] (p)--node[above left]{80\%} (m1); \draw[->] (p)--node[above right]{20\%} (m2);
\draw[->] (m1)--(c1); \draw[->] (m2)--(c2); \draw[->] (c1)--(f1); \draw[->] (c2)--(f2);
\end{tikzpicture}
\caption{Обмен строками между процессами}
\end{figure}
\subsection{Организация обмена}
Объект памяти создаётся через \texttt{shm\_open()} и \texttt{ftruncate()}, затем отображается вызовом \texttt{mmap()} с флагом \texttt{MAP\_SHARED}. В каждом буфере есть массив данных на 4096 байт, длина фрагмента, признаки последнего фрагмента и завершения передачи, а также два межпроцессных семафора. Семафор \texttt{empty} сначала равен единице, \texttt{full} -- нулю. Записывающий процесс ждёт освобождения буфера, помещает фрагмент и сигнализирует о готовности; читающий ждёт данные, копирует их и освобождает буфер. Строки длиннее 4096 байт передаются по частям и собираются в дочернем процессе.

После конца ввода родитель отправляет каждому процессу сообщение с признаком завершения и вызывает \texttt{waitpid()}. Затем отображения и дескрипторы закрываются, а владелец удаляет объекты памяти и семафоры. В коде проверяются результаты основных системных вызовов; ожидание семафоров и процессов повторяется при прерывании сигналом \texttt{EINTR}.
\section{Ключевые фрагменты кода}
Ниже приведены фрагменты из исходных файлов проекта. Полные файлы доступны в репозитории.
\clearpage
\subsection{Создание отображения и семафоров}
Объект общей памяти получает размер структуры буфера, после чего в нём создаются семафоры, доступные разным процессам.
\begin{lstlisting}
new_memory->fd = shm_open(name, O_CREAT | O_EXCL | O_RDWR, 0600);
ftruncate(new_memory->fd, sizeof(shared_buffer));
void *mapping = mmap(NULL, sizeof(shared_buffer),
    PROT_READ | PROT_WRITE, MAP_SHARED, new_memory->fd, 0);
new_memory->buffer = mapping;
sem_init(&new_memory->buffer->empty, 1, 1);
sem_init(&new_memory->buffer->full, 1, 0);
\end{lstlisting}
\begin{center}\small Листинг 1 -- Инициализация общей памяти (проверки ошибок опущены)\end{center}
\subsection{Передача фрагмента строки}
Отправитель дожидается свободного буфера и сообщает получателю о появлении данных. Признак \texttt{is\_last} позволяет получателю определить конец строки.
\begin{lstlisting}
size_t part = length - offset;
if (part > BUFFER_SIZE) {
    part = BUFFER_SIZE;
}
if (wait_semaphore(&buffer->empty) == -1) {
    perror("sem_wait(empty)");
    return -1;
}
buffer->length = part;
buffer->is_last = (offset + part == length);
buffer->finished = 0;
if (part > 0) {
    memcpy(buffer->data, data + offset, part);
}
if (sem_post(&buffer->full) == -1) {
    perror("sem_post(full)");
    return -1;
}
offset += part;
\end{lstlisting}
\begin{center}\small Листинг 2 -- Запись одного фрагмента строки\end{center}
\clearpage
\subsection{Распределение и обработка строк}
В родительском процессе выбор получателя выполняется для каждой прочитанной строки.
\begin{lstlisting}
int choice = rand() % 100;
if (choice < 80) {
    os_shm_write(memory1, line, length);
} else {
    os_shm_write(memory2, line, length);
}
\end{lstlisting}
\begin{center}\small Листинг 3 -- Выбор дочернего процесса (проверки ошибок опущены)\end{center}

Дочерняя программа читает строку из своего объекта памяти, удаляет гласные функцией \texttt{remove\_vowels()} и пишет оставшиеся символы в файл через стандартный вывод.
\begin{lstlisting}
int result = os_shm_read(memory, &line, &length);
remove_vowels(line);
if (fputs(line, stdout) == EOF) {
    free(line);
    os_shm_destroy(memory);
    return 1;
}
free(line);
\end{lstlisting}
\begin{center}\small Листинг 4 -- Обработка строки дочерним процессом\end{center}
\section{Состояние реализации и проверка}
При анализе исходных файлов обнаружены несогласованные объявления интерфейса: в \texttt{os\_api.h} третий аргумент \texttt{os\_shm\_read} объявлен как \texttt{size\_t}, тогда как реализация и вызов используют \texttt{size\_t *}; второй аргумент \texttt{os\_shm\_write} объявлен как \texttt{const void *}, а в определении -- как \texttt{const char *}. Кроме того, заголовок объявляет \texttt{os\_destroy\_shm}, тогда как реализация и вызывающий код используют \texttt{os\_shm\_destroy}. Эти расхождения необходимо исправить перед сборкой.

Проверка запуска и получение журналов процессов не выполнены: доступная среда WSL вернула системную ошибку при запуске. Поэтому фактический консольный вывод, идентификаторы процессов и содержимое результирующих файлов в отчёте не приводятся. В проекте также нет сценария сборки или сохранённых результатов тестов. Распределение 80/20 описывает правило выбора для отдельных строк; точное соотношение в конечной выборке не гарантируется. При аварийном завершении дочернего процесса ожидание семафора родителем может остаться заблокированным.
\section{Выводы}
В проекте подготовлена реализация обмена строками между родительским и двумя дочерними процессами через отображаемую POSIX-память. Передача строк разделена на фрагменты, доступ к буферам синхронизируется семафорами, а результат дочерних процессов направляется в отдельные файлы. В ходе работы изучены принципы отображения памяти с помощью вызова \texttt{mmap()}, синхронизации процессов семафорами и передачи данных через общий буфер.
\end{document}









